% PAGES: 181-182
% Closes the "Proof of Theorem 6.1" opened in sec06_c1_1.tex (page 178 /
% folio 162) and continued through sec06_c1_2.tex (page 179-180 / folio
% 163-164): the "\end{proof}" a few lines below matches that
% "\begin{proof}[Proof of Theorem 6.1]". No other environment is open on
% entry.

Then, from (6.15) and (6.18), we obtain that
\begin{equation}
(U(2:n,2:n))\tp U(2:n,2:n) \;=\; A(2:n,2:n) \;-\; \frac{(A(1,2:n))\tp A(1,2:n)}{A(1,1)}.
\end{equation}

The crucial step that makes the Cholesky decomposition possible is that, if $A$ is a
symmetric positive definite matrix, then the matrix
\begin{equation}
A(2:n,2:n) \;-\; \frac{(A(1,2:n))\tp A(1,2:n)}{A(1,1)}
\end{equation}
from the right hand side of (6.19) is symmetric positive definite matrix. This result
is proved in Chapter 10; see Lemma 10.17.

Then, if follows from (6.19) that the $(n-1) \times (n-1)$ matrix $U(2:n,2:n)$ is the
Cholesky factor of the symmetric positive definite matrix from (6.20), and we conclude
that the matrix $A$ has a Cholesky decomposition $A = U\tp U$ with Cholesky factor
\[
U \;=\; \begin{pmatrix} U(1,1) & U(1,2:n) \\ 0 & U(2:n,2:n) \end{pmatrix}.
\]
\end{proof}

The proof above also provides the idea for the recursive implementation of the
Cholesky decomposition algorithm. Once the first row of $U$ is computed, the entries of
$A(2:n,2:n)$ are updated\footnote{The entries of $A(2:n,2:n)$ are updated using
(6.15), i.e., $A(2:n,2:n) - (U(1,2:n))\tp U(1,2:n)$; formula (6.19), i.e., $A(2:n,2:n) -
\frac{(A(1,2:n))\tp A(1,2:n)}{A(1,1)}$, is only used for the proof of the existence of
the Cholesky decomposition.} to $A(2:n,2:n) - (U(1,2:n))\tp U(1,2:n)$, i.e.,
\begin{equation}
A(2:n,2:n) \;=\; A(2:n,2:n) \;-\; (U(1,2:n))\tp U(1,2:n).
\end{equation}
Then, the first row of the matrix $U(2:n,2:n)$ (which coincides with the second row of
$U$ without the first entry which is equal to $0$) is computed from the matrix
$A(2:n,2:n) - (U(1,2:n))\tp U(1,2:n)$ in the same way as the first row of $U$ was
computed from the matrix $A$. The matrix $A(3:n,3:n)$ is then updated and will again be
a symmetric positive definite matrix, and therefore this recursive process continues
until all the rows of $U$ are computed.

The Cholesky factor $U$ is obtained once all its $n$ rows are computed recursively via
this algorithm; see section 6.1.1 for more details.

\begin{theorem}
The Cholesky decomposition of a symmetric positive definite matrix is unique.
\end{theorem}

\begin{proof}
We give a proof by contradiction. Assume that the symmetric positive definite matrix
$A$ of size $n$ has two Cholesky decompositions, i.e., assume that
\[
A \;=\; U_1\tp U_1 \quad \text{and} \quad A \;=\; U_2\tp U_2,
\]
where $U_1$ and $U_2$ are upper triangular matrices with positive entries on the main
diagonal. Then,
\begin{equation}
U_1\tp U_1 \;=\; U_2\tp U_2.
\end{equation}
Note that $U_1$ and $U_2$ are nonsingular matrices, since any Cholesky factor is a
nonsingular matrix; see Definition 6.1. Multiply (6.22) by $(U_2\tp)^{-1}$ on the left
and by $U_1^{-1}$ on the right and obtain
\begin{align}
(U_2\tp)^{-1}(U_1\tp U_1)U_1^{-1} &= (U_2\tp)^{-1}(U_2\tp U_2)U_1^{-1} \notag \\
\Longleftrightarrow \quad \big((U_2\tp)^{-1}U_1\tp\big) \cdot \big(U_1 U_1^{-1}\big)
&= \big((U_2\tp)^{-1}U_2\tp\big) \cdot \big(U_2 U_1^{-1}\big) \notag \\
\Longleftrightarrow \quad (U_2\tp)^{-1}U_1\tp &= U_2 U_1^{-1},
\end{align}
since $U_1U_1^{-1} = I$ and $(U_2\tp)^{-1}U_2\tp = I$.

Recall from Lemma 1.17 that the inverse of an upper triangular matrix is upper
triangular, and the inverse of a lower triangular matrix is lower triangular, and, from
Lemma 1.15, that the product of two upper triangular matrices is upper triangular, and
the product of two lower triangular matrices is lower triangular. Thus, the matrix
$U_1^{-1}$ is upper triangular and therefore the matrix $U_2U_1^{-1}$ is also upper
triangular. Similarly, the matrix $(U_2\tp)^{-1}$ is lower triangular and therefore the
matrix $(U_2\tp)^{-1}U_1\tp$ is also lower triangular. Since the matrices
$(U_2\tp)^{-1}U_1\tp$ and $U_2U_1^{-1}$ are equal, see (6.23), it follows that they
must be diagonal matrices.

Let $D = \diag(d_j)_{j=1:n}$ be a diagonal matrix such that
\[
D \;=\; (U_2\tp)^{-1}U_1\tp \;=\; U_2U_1^{-1}.
\]
Then, by multiplying $D = U_2U_1^{-1}$ to the right by $U_1$, we find that
\begin{equation}
U_2 \;=\; DU_1
\end{equation}
and therefore
\begin{equation}
A \;=\; U_2\tp U_2 \;=\; (DU_1)\tp U_2 \;=\; U_1\tp D\tp DU_1 \;=\; U_1\tp D^2 U_1,
\end{equation}
where the last equality follows from the fact that $D\tp = D$, since $D$ is a diagonal
matrix. Since $A = U_1\tp U_1$ and $U_1$ is a nonsingular matrix, it follows from
(6.25) that
\begin{align*}
U_1\tp U_1 = U_1\tp D^2 U_1 &\iff (U_1\tp)^{-1}U_1\tp U_1(U_1)^{-1} = (U_1\tp)^{-1}U_1\tp
D^2 U_1(U_1)^{-1} \\
&\iff I = D^2.
\end{align*}
Note that $D^2 = \diag(d_j^2)_{j=1:n}$, see (1.93), and therefore $I = D^2$ if and only
if $d_j^2 = 1$ for all $j = 1:n$. Thus, all the diagonal entries of $D$ are either $1$
or $-1$, i.e.,
\begin{equation}
d_j = 1 \quad \text{or} \quad d_j = -1, \quad \forall\, j = 1:n.
\end{equation}
Recall from (6.24) that $U_2 = DU_1$. From Lemma 1.10, we obtain that the $j$-th row of
the matrix $DU_1$ is equal to the $j$-th row of the matrix $U_1$ multiplied by $d_j$.
In particular, $(DU_1)(j,j) = d_jU_1(j,j)$. Since $U_2 = DU_1$, it follows that
$U_2(j,j) = d_jU_1(j,j)$, and therefore
\begin{equation}
d_j \;=\; \frac{U_2(j,j)}{U_1(j,j)} \;>\; 0, \quad \forall\, j = 1:n,
\end{equation}
since, by Definition 6.1, all the main diagonal entries of $U_1$ and $U_2$ are
positive.

From (6.26) and (6.27), we obtain that $d_j = 1$ for all $j = 1:n$, and therefore the
matrix $D = \diag(d_j)_{j=1:n}$ is equal to the identity matrix, i.e., $D = I$.

Since $D = I$, it follows from (6.24) that $U_2 = U_1$, and we conclude that the
Cholesky decomposition of the matrix $A$ is unique.
\end{proof}

\subsection{Pseudocode and operation count for Cholesky decomposition}

Based on the proof of Theorem 6.1, the algorithm for finding the Cholesky decomposition
of a symmetric positive definite matrix can be implemented recursively as detailed
below; see also the pseudocode from Table 6.1.

Let $A$ be a symmetric positive definite matrix and let $U$ be the Cholesky factor of
$A$.

\begin{itemize}
\item Compute the first row of $U$:
\end{itemize}
